% ============================================================
% DOCUMENTCLASS
% ============================================================
\documentclass[12pt,oneside]{book}

% ============================================================
% METADATOS
% ============================================================
\newcommand{\universidad}{Universidad de Buenos Aires (UBA)}
\newcommand{\facultad}{Facultad de Derecho}
\newcommand{\materia}{Derecho Procesal Civil}
\newcommand{\titulo}{La Prueba Documental y la Prueba Informativa (Unidad 13)}
\newcommand{\autor}{Isaac Edgar Camacho Ocampo}
\newcommand{\anio}{2024 - 2025}

% ============================================================
% PAQUETES
% ============================================================
\usepackage[utf8]{inputenc}
\usepackage[T1]{fontenc}
\usepackage[spanish,es-nodecimaldot]{babel}

\usepackage[
    top=2.5cm,
    bottom=2.5cm,
    left=3cm,
    right=2.5cm,
    headheight=15pt
]{geometry}

\usepackage{amsmath}
\usepackage{amssymb}
\usepackage{mathtools}

\usepackage{graphicx}
\usepackage{float}
\usepackage{caption}
\usepackage{subcaption}

\usepackage{booktabs}
\usepackage{array}
\usepackage{longtable}
\usepackage{tabularx}

\usepackage{enumitem}

\usepackage{xcolor}
\usepackage[most]{tcolorbox}

\usepackage{fancyhdr}
\usepackage{ragged2e}

% ============================================================
% CONFIGURACIÓN
% ============================================================
\graphicspath{
    {./img/}
    {./graficos/}
    {./figuras/}
}

\setcounter{tocdepth}{2}

\setlength{\parindent}{0pt}
\setlength{\parskip}{0.5em}

\setlist[itemize]{
    topsep=0.3em,
    itemsep=0.2em
}

\setlist[enumerate]{
    topsep=0.3em,
    itemsep=0.2em
}

\clubpenalty=10000
\widowpenalty=10000

% ============================================================
% COMANDOS
% ============================================================
\newcommand{\abs}[1]{\left|#1\right|}
\newcommand{\norm}[1]{\left\lVert#1\right\rVert}
\newcommand{\vect}[1]{\mathbf{#1}}
\newcommand{\deriv}[2]{\frac{d#1}{d#2}}
\newcommand{\pderiv}[2]{\frac{\partial#1}{\partial#2}}

% ============================================================
% ENTORNOS / CAJAS
% ============================================================
\newtcolorbox{definition}[1]{
    enhanced,
    breakable,
    title={Definición: #1},
    fonttitle=\bfseries,
    colback=blue!3,
    colframe=blue!50!black,
    boxrule=0.8pt,
    arc=2mm
}

\newtcolorbox{important}{
    enhanced,
    breakable,
    title={Importante},
    fonttitle=\bfseries,
    colback=yellow!8,
    colframe=orange!70!black,
    boxrule=0.8pt,
    arc=2mm
}

\newtcolorbox{example}[1]{
    enhanced,
    breakable,
    title={Ejemplo: #1},
    fonttitle=\bfseries,
    colback=green!4,
    colframe=green!40!black,
    boxrule=0.8pt,
    arc=2mm
}

\newtcolorbox{formula}[1]{
    enhanced,
    breakable,
    title={Fórmula / Regla: #1},
    fonttitle=\bfseries,
    colback=purple!4,
    colframe=purple!50!black,
    boxrule=0.8pt,
    arc=2mm
}

\newtcolorbox{exercise}[1]{
    enhanced,
    breakable,
    title={Ejercicio: #1},
    fonttitle=\bfseries,
    colback=gray!5,
    colframe=gray!60!black,
    boxrule=0.8pt,
    arc=2mm
}

\newtcolorbox{note}{
    enhanced,
    breakable,
    title={Nota},
    fonttitle=\bfseries,
    colback=white,
    colframe=gray!50,
    boxrule=0.6pt,
    arc=2mm
}

% ============================================================
% CONFIGURACIÓN FINAL (encabezados, hipervínculos)
% ============================================================
\pagestyle{fancy}
\fancyhf{}

\fancyhead[L]{\nouppercase{\leftmark}}
\fancyhead[R]{\materia}

\fancyfoot[C]{\thepage}

\renewcommand{\headrulewidth}{0.4pt}
\renewcommand{\footrulewidth}{0pt}

\fancypagestyle{plain}{
    \fancyhf{}
    \fancyfoot[C]{\thepage}
    \renewcommand{\headrulewidth}{0pt}
}

\usepackage[
    colorlinks=true,
    linkcolor=blue!50!black,
    citecolor=green!40!black,
    urlcolor=blue!60!black,
    pdftitle={\titulo},
    pdfauthor={\autor},
    pdfsubject={Apuntes universitarios},
    bookmarks=true
]{hyperref}

% ============================================================
% DOCUMENT
% ============================================================
\begin{document}

% ------------------------------------------------------------
% PORTADA
% ------------------------------------------------------------
\begin{titlepage}

\thispagestyle{empty}

\begin{center}

\vspace*{1cm}

{\large \textsc{\universidad}}

\vspace{0.2cm}

{\large \textsc{\facultad}}

\vspace{3cm}

{\Huge \textbf{\materia}}

\vspace{1cm}

{\LARGE \titulo}

\vspace{0.8cm}

\textit{Comprender los conceptos es el primer paso para convertir el estudio en conocimiento.}

\vspace{1.2cm}

\rule{0.70\textwidth}{0.5pt}

\vspace{0.5cm}

{\large Apuntes de estudio}

\vspace{0.5cm}

\rule{0.70\textwidth}{0.5pt}

\vfill

{\large \autor}

\vspace{0.3cm}

{\large \anio}

\end{center}

\vfill

\begin{flushleft}
\textit{Este es un modesto aporte a los estudiantes de ``Derecho Procesal Civil'' no pretende sustituir el material oficial de la materia.}
\end{flushleft}

\end{titlepage}

% ------------------------------------------------------------
% ÍNDICE
% ------------------------------------------------------------
\tableofcontents

% ------------------------------------------------------------
% PRESENTACIÓN GENERAL
% ------------------------------------------------------------
\chapter*{Presentación}
\addcontentsline{toc}{chapter}{Presentación}

Como todo contrato o acuerdo de la vida cotidiana necesita de un ``papel que lo respalde'', el proceso judicial necesita de la prueba documental para demostrar los hechos. El presente material recorre desde el ticket de estacionamiento hasta la escritura pública, y desde el correo electrónico hasta la firma digital, mostrando cómo el derecho regula cada uno de esos ``papeles'' que se usan a diario. Asimismo, se estudia la prueba informativa como medio para incorporar al proceso datos registrados por terceros.

Los medios de prueba abordados en esta unidad se enmarcan dentro de los medios de prueba previstos por el Código Procesal, que tipifica expresamente: la prueba documental, la informativa, la confesional, la testimonial, la pericial y el reconocimiento judicial.

\section*{Esquema general del contenido}
\addcontentsline{toc}{section}{Esquema general del contenido}

A continuación se presenta un esquema orientativo de los bloques temáticos desarrollados en este material.

\begin{longtable}{p{0.08\textwidth}p{0.87\textwidth}}
\toprule
\multicolumn{2}{l}{\textbf{Bloque I: Concepto y clasificación de documento}} \\
\midrule
1.1 & Concepto amplio de documento: elemento, pensamiento, soporte, tangible, digital. \\
1.2 & Clasificación general: papel vs. electrónico, firmado vs. no firmado. \\
1.3 & Impacto del nuevo Código Civil 2015: aggiornamiento, instrumentos, regulación moderna. \\
\midrule
\multicolumn{2}{l}{\textbf{Bloque II: Documentos en soporte papel}} \\
\midrule
2.1 & Instrumentos públicos (art. 289 CCyC): escritura, oficial público, títulos estatales. \\
2.2 & Valor probatorio del instrumento público (art. 296 CCyC): plena fe, lugar, fecha, acto celebrado. \\
2.3 & Instrumentos privados: firma, doble ejemplar, reconocimiento, cotejo. \\
2.4 & Autenticación y pericia caligráfica (arts. 390/458 CPCC): dubitado, indubitado, cuerpo de escritura. \\
2.5 & Abuso de firma en blanco (art. 315 CCyC): principio de prueba escrito, contradocumento. \\
2.6 & Instrumentos particulares no firmados (art. 319 CCyC): ticket, voucher, valor a criterio judicial. \\
\midrule
\multicolumn{2}{l}{\textbf{Bloque III: Documentos electrónicos y correspondencia}} \\
\midrule
3.1 & Firma digital (Ley 25.506): entidad certificante, plena fe, impugnación técnica. \\
3.2 & Firma electrónica: homebanking, clave personal, carga de autenticación. \\
3.3 & Documentos electrónicos no firmados: videos, audios, Street View, valoración judicial. \\
3.4 & Correspondencia (arts. 318/319 CCyC): epistolar, email, confidencial, remitente. \\
3.5 & Libros contables (art. 330 CCyC): valor contra quien los lleva, indivisibilidad. \\
\midrule
\multicolumn{2}{l}{\textbf{Bloque IV: Ofrecimiento e impugnación de prueba documental}} \\
\midrule
4.1 & Oportunidad del ofrecimiento (arts. 333/335 CPCC): demanda, contestación, documentos posteriores. \\
4.2 & Documentos en poder de la contraparte (art. 388 CPCC): intimación, presunción contraria, negativa. \\
4.3 & Impugnación del instrumento público (art. 395 CPCC): redargución de falsedad, incidente, 10 días. \\
4.4 & Impugnación del instrumento privado: cotejo, pericia caligráfica, autenticación judicial. \\
\midrule
\multicolumn{2}{l}{\textbf{Bloque V: Prueba informativa}} \\
\midrule
5.1 & Concepto y diferencia con prueba testimonial: datos registrados, terceros, archivos. \\
5.2 & Formato del oficio (arts. 398/400 CPCC): estructura, transcripción del auto, firma del abogado. \\
5.3 & Plazo y caducidad (art. 402 CPCC): 10 días entidad, 5 días reiteración, caducidad. \\
5.4 & Sistemas electrónicos modernos: digitalización, bancos, migraciones, registros. \\
\bottomrule
\end{longtable}

\begin{important}
Todo manual o tratado que no contemple la reforma del Código Civil de 2015 estará desactualizado en materia de prueba documental.
\end{important}

% ============================================================
% CAPÍTULO 1
% ============================================================
\chapter{Concepto de documento}

Dentro de los medios de prueba previstos por el Código Procesal se encuentran expresamente tipificados: la prueba documental, la informativa, la confesional, la testimonial, la pericial y el reconocimiento judicial.

\begin{definition}{Documento en sentido lato}
En el sentido lato, documento es todo elemento susceptible de representar una manifestación del pensamiento, con prescindencia de la forma en que esté plasmada o registrada esa representación.
\end{definition}

Tradicionalmente el documento estaba más asociado a un objeto tangible, corpóreo, preponderantemente extendido por escrito. Sin embargo, la realidad actual exige ampliar significativamente ese concepto.

\section{El documento en la era digital}

Hoy en día, dentro de la idea de documento, se encuentran registros que pertenecen a algo poco tangible porque está contenido en información archivada en formato binario. Se trata de registros informáticos fácilmente rescatables que no están necesariamente en un soporte material: el soporte puede estar en la nube o en algún tipo de soporte magnético.

Por tanto, \textbf{documento es todo elemento susceptible de representar una manifestación del pensamiento humano, independientemente de su soporte}.

\section{Impacto del Código Civil y Comercial (2015)}

El Código Civil y Comercial, sancionado y en vigencia desde el año 2015, vino a ``aggiornar'' el derecho en lo relativo a los documentos. Incorporó material procesal con regulación de instrumentos públicos, instrumentos privados, instrumentos particulares y los documentos digitales o electrónicos.

\begin{important}
El Código Procesal, pensado desde la perspectiva del viejo Código Civil, no tuvo ese aggiornamiento. Por eso, al estudiar prueba documental, es fundamental recurrir a textos y manuales posteriores a 2015.
\end{important}

% ============================================================
% CAPÍTULO 2
% ============================================================
\chapter{Clasificación de los documentos}

La clasificación principal de los documentos distingue entre:

\begin{enumerate}[label=\alph*)]
    \item \textbf{Documentos en soporte papel.} Con sus subcategorías de instrumentos públicos e instrumentos particulares.
    \item \textbf{Documentos en soporte electrónico.} Con las subcategorías de firma digital, firma electrónica y documentos sin firma.
\end{enumerate}

El Código Civil y Comercial también regula expresamente la correspondencia (artículos 318 y 319) y los libros de comercio, ahora contenidos en el CCyC desde la derogación del Código de Comercio.

% ============================================================
% CAPÍTULO 3
% ============================================================
\chapter{Documentos en soporte papel}

\section{Instrumentos públicos}

El artículo 289 del Código Civil y Comercial divide los instrumentos públicos en tres grandes grupos:

\begin{enumerate}
    \item Las escrituras públicas.
    \item Los instrumentos extendidos por un oficial público de acuerdo a las normas establecidas en cuanto a las formas por la ley.
    \item Los títulos emitidos por los estados nacional, provincial o municipal, de acuerdo a la legislación correspondiente.
\end{enumerate}

\begin{formula}{Artículo 289 CCyC — Instrumentos públicos}
Son instrumentos públicos: (a) las escrituras públicas y sus copias o testimonios; (b) los instrumentos que extienden los escribanos o funcionarios públicos con los requisitos que establecen las leyes; (c) los títulos emitidos por el Estado nacional, provincial o la Ciudad Autónoma de Buenos Aires, conforme a las leyes que autorizan su emisión.
\end{formula}

\subsection{Valor probatorio del instrumento público}

Tradicionalmente —y el nuevo código lo recepta— se sostiene que estos documentos hacen plena fe hasta tanto no haya una decisión judicial que les quite toda validez a través de un proceso de impugnación. En principio hacen plena fe hasta que no exista una declaración de que resultan falsos.

El artículo 296 del Código Civil y Comercial fue específico en determinar qué aspectos del instrumento público gozan de esta fe pública.

\begin{formula}{Artículo 296 CCyC — Eficacia probatoria}
El instrumento público hace plena fe: (a) en cuanto a que se ha realizado el acto, la fecha, el lugar y los hechos que el oficial público enuncia como cumplidos por él o ante él, hasta que sea declarado falso en juicio civil o criminal; (b) en cuanto al contenido de las declaraciones sobre convenciones, disposiciones, pagos, reconocimientos y enunciaciones de hechos directamente relacionados con el objeto principal del acto instrumentado, hasta que se produzca prueba en contrario.
\end{formula}

Es posible diferenciar claramente dos aspectos del instrumento público:

\begin{example}{El escribano y la manifestación de las partes}
Si el escribano dice que frente a él las partes manifestaron que antes de este momento se entregaron una suma de dinero, pero esa entrega no sucede frente al escribano, esa circunstancia vale como manifestación de las partes frente al escribano. El escribano da plena fe de que las partes dicen que antes de este momento se hizo entrega material de la posesión del inmueble, por ejemplo; pero no hace plena fe de que esa entrega de posesión fue real.
\end{example}

\begin{table}[H]
\centering
\caption{Aspectos del instrumento público y su forma de impugnación}
\begin{tabularx}{\textwidth}{>{\raggedright\arraybackslash}X >{\raggedright\arraybackslash}p{0.28\textwidth} >{\raggedright\arraybackslash}X}
\toprule
\textbf{Aspecto del instrumento público} & \textbf{Fe pública} & \textbf{Forma de impugnar} \\
\midrule
Celebración del acto, lugar, fecha y hechos ocurridos ante el oficial público & Plena fe -- irrefutable hasta sentencia & Solo por redargución de falsedad (incidente o acción penal/civil) \\
Contenido de las declaraciones de las partes ante el escribano (pagos reconocidos, enunciaciones) & No hace plena fe -- admite prueba en contrario & Cualquier medio de prueba en el juicio principal \\
\bottomrule
\end{tabularx}
\end{table}

\section{Instrumentos privados}

El instrumento privado es el instrumento particular firmado. Es el que goza de la formalidad que el Código establece, fundamentalmente la firma de las partes.

\subsection{La firma en el nuevo Código Civil}

El nuevo Código Civil incorporó al concepto de firma una serie de datos que ya estaban en la jurisprudencia y la doctrina, pero no en el Código de Vélez:

\begin{itemize}
    \item La firma no necesariamente debe ser escritura del nombre de la persona; pueden ser los signos habituales que la persona realiza como marca distintiva a la hora de suscribir un documento.
    \item Se equipara a la firma la impresión dígito-pulgar.
    \item Se regula la ``firma a ruego'': cuando otra persona estampa la firma a pedido de alguien que no sabe o no puede escribir.
\end{itemize}

\subsection{El doble ejemplar}

En los actos jurídicos bilaterales debe haber, para cada interesado, un ejemplar que le permita hacer valer frente a los otros participantes el contenido del acto. No se requieren necesariamente dos ejemplares: deben expedirse tantos ejemplares como partes hayan intervenido en el acto.

\subsection{Valor probatorio del instrumento privado}

A diferencia de la escritura pública, que vale por sí misma en virtud de la intervención del oficial público, el valor probatorio del instrumento privado es sólido en la medida en que se encuentra reconocido o autenticado.

\begin{formula}{Artículo 356 CPCC — Cargas del demandado}
El demandado deberá, al contestar la demanda: 1) Reconocer o negar categóricamente cada uno de los hechos expuestos en la demanda... 2) Reconocer o negar la autenticidad de los documentos acompañados que se le atribuyen y la recepción de las cartas y telegramas a él dirigidos cuyas copias se acompañen... El silencio, las respuestas evasivas o la negativa meramente general podrán estimarse como reconocimiento de la verdad de los hechos pertinentes y lícitos a que se refieran.
\end{formula}

El incumplimiento de esa carga importa tenerlos por reconocidos, tanto los instrumentos privados firmados como las cartas atribuidas al demandado.

\begin{important}
Con el nuevo Código, el valor probatorio del instrumento privado reconocido puede ser incluso mayor que el del instrumento público, dado que el único modo de impugnarlo es plantear la nulidad por vicios del consentimiento a la hora de reconocer ese instrumento.
\end{important}

\subsection{El cotejo y la pericia caligráfica}

Cuando el instrumento privado presentado en juicio es negado por quien supuestamente lo firmó, es necesario realizar el cotejo de esa firma con firmas auténticas de la persona que niega que la firma le pertenezca. Este cotejo se realiza a través de la prueba pericial caligráfica, regulada en el artículo 390 del CPCC, que remite al artículo 458 y siguientes.

\begin{formula}{Artículo 390 CPCC — Reconocimiento o cotejo}
Cuando se niegue la autenticidad de un documento privado, o se alegue falsedad o adulteración, la parte que lo hubiere presentado podrá pedir el reconocimiento por la persona a quien se atribuya o el cotejo, si fuere posible.
\end{formula}

El perito calígrafo necesitará documentos indubitados y dubitados:

\begin{itemize}
    \item \textbf{Documentos indubitados}: aquellos que indudablemente pertenecen al puño y letra de la persona a quien se le atribuye la firma. Son la base del cotejo.
    \item \textbf{Documento dubitado}: el documento cuya firma está en duda y debe ser cotejada.
\end{itemize}

Si las partes no se ponen de acuerdo sobre cuáles son los documentos indubitados, el juez puede disponer la elaboración de un cuerpo de escritura: se cita a la persona cuya firma se cuestiona y, en presencia del perito calígrafo, el secretario del juzgado y las partes, se le pide que escriba y firme una serie de palabras que se le dictan en una hoja en blanco, obteniendo así un cuerpo de escritura indubitado.

\subsection{El caso del fallecido y registros alternativos}

Cuando la persona cuya firma se cuestiona ha fallecido (caso frecuente en la impugnación de testamentos ológrafos), se puede recurrir a distintos registros donde puedan obrar firmas indubitadas:

\begin{itemize}
    \item Registros de cuentas bancarias (archivo de firmas de titulares).
    \item Registro Nacional de las Personas.
    \item Escrituras públicas (protocolo notarial).
    \item Registros policiales (en épocas anteriores, la Policía Federal llevaba legajo de antecedentes).
\end{itemize}

\subsection{El abuso de firma en blanco (art. 315 CCyC)}

El artículo 315 del Código Civil prevé la posibilidad de invocar que existió un abuso de firma en blanco: cuando una persona suscribe una hoja en blanco por razones de confianza o manejo de ciertos negocios, y la otra parte es acusada de haber abusado de esa hoja en blanco, llenándola con un contenido distinto al convenido.

\begin{formula}{Artículo 315 CCyC — Abuso de firma en blanco}
Para impugnar un instrumento al que se atribuye un abuso de firma en blanco se requiere un principio de prueba por escrito. Se puede presentar todo tipo de prueba para sostener la impugnación, pero necesariamente debe acompañarse, además de la prueba testimonial, un principio de prueba por escrito. No bastan testimonios solos.
\end{formula}

La existencia de un contradocumento —otro documento escrito que contraviene lo que dice el instrumento anterior— puede ser un elemento sustancial a la hora de este tipo de demostraciones, aunque no alcanza por sí solo.

\subsection{Instrumento particular con impresión dígito-pulgar}

El nuevo Código también contempla el instrumento particular firmado con impresión dígito-pulgar. Sin embargo, no le da la misma validez que una rúbrica con firma: es considerado un principio de prueba que puede ser impugnado por cualquier medio.

\begin{note}
Quien firma con impresión dígito-pulgar generalmente es una persona que no sabe leer y escribir, lo que facilita que se le induzca a suscribir algo distinto de lo que se le dice. Por eso, para mayor seguridad, es preferible siempre recurrir a la ``firma a ruego''.
\end{note}

\section{Instrumentos particulares (sin firma)}

Los instrumentos particulares son los extendidos por escrito sin constancia de firma. A esta categoría pertenecen infinidad de actos y contratos cotidianos.

\begin{example}{El comprobante de estacionamiento}
Cuando uno deja el auto en una playa de estacionamiento, como único comprobante de que ha dejado en esa plaza un auto y las llaves, el playero le extiende un papelito sin firma alguna que solo dice los tres últimos números de la patente. Esa constancia --eso no es ni más ni menos que un instrumento particular-- sin firma alguna, ¡y vale como elemento de prueba del depósito o contrato de garaje!
\end{example}

Otros ejemplos: vouchers de viaje, tickets, cupones, boletos. El contrato de transporte del boleto de colectivo hoy ya no es un instrumento particular porque es un registro electrónico, pero cuando las máquinas extendían esa impresión en papel, ese instrumento particular era una constancia demostrativa del contrato de transporte.

El valor probatorio de estos instrumentos queda a discreción del juez, quien debe evaluar la congruencia entre lo sucedido, lo narrado y la calidad técnica del documento. Algunos instrumentos particulares pueden incorporar códigos de barras u otros elementos técnicos que requieren interpretación técnica para descifrar lo que surge de ellos, lo que puede darles mayor confiabilidad.

\begin{formula}{Artículo 319 CCyC — Valor probatorio de los instrumentos particulares}
El valor probatorio de los instrumentos particulares debe ser apreciado por el juez ponderando, entre otras pautas, la congruencia entre lo sucedido y narrado, la precisión y claridad técnica del texto, los usos y prácticas del tráfico, las relaciones precedentes y la confiabilidad de los soportes utilizados y de los procedimientos técnicos que se apliquen.
\end{formula}

% ============================================================
% CAPÍTULO 4
% ============================================================
\chapter{Documentos electrónicos}

Los documentos electrónicos son todos aquellos plasmados en un soporte informático. La hipertrofia de este tipo de documentos en negociaciones online hace necesaria la regulación que el Código Civil ha introducido.

Se distinguen tres tipos, paralelos a los documentos en soporte papel: firma digital, firma electrónica y documentos electrónicos sin firma.

\section{Firma digital (Ley 25.506)}

La firma digital está regulada por la Ley 25.506. Se da en los supuestos en donde existe un soporte electrónico que posee una clave oficialmente certificada por un organismo del Estado. A través de un mecanismo de algoritmos, se identifica que el emisor de esa firma ha concurrido previamente a la entidad pública certificante y ha obtenido la clave que le permite operar.

\begin{formula}{Ley 25.506 — Firma Digital}
La firma digital es el resultado de aplicar a un documento digital un procedimiento matemático que requiere información de exclusivo conocimiento del firmante, encontrándose ésta bajo su absoluto control. La firma digital debe ser susceptible de verificación por terceras partes, de manera tal que dicha verificación simultáneamente permita identificar al firmante y detectar cualquier alteración del documento digital posterior a su firma.
\end{formula}

El organismo estatal certificante garantiza la identidad de quienes suscriben y la autenticidad del instrumento. Por eso los instrumentos con firma digital hacen plena fe en el proceso, tanto de la autenticidad de los contenidos como de la identidad del emisor.

% TODO: contenido incompleto o ambiguo en las notas originales — se afirma que la firma digital posee incluso menor posibilidad de falibilidad que los instrumentos públicos pasados ante oficial público, sin precisarse la fuente de esa afirmación dentro del apunte.
Se sostiene que la firma digital posee incluso menor posibilidad de falibilidad que los instrumentos públicos pasados ante oficial público, dado que el sistema de encriptamiento es técnicamente muy confiable.

Para impugnar una firma digital, quien la cuestiona tiene la carga de demostrar técnicamente —mediante pericia informática— que existió algún defecto en el sistema de certificación, en la emisión o en la vigencia de la firma.

\section{Firma electrónica}

La firma electrónica no es firma digital. Es la que se utiliza habitualmente en:

\begin{itemize}
    \item Homebanking.
    \item MercadoPago y compras online.
    \item Compras con tarjeta de crédito.
    \item Presentaciones judiciales en el sistema de gestión judicial.
\end{itemize}

No existe una entidad pública que esté certificando los clics que uno hace en un homebanking o en una compra online. Lo que hay son claves que permiten, cruzando datos, identificar al sujeto que está prestando conformidad con determinadas operaciones.

\begin{example}{La clave personal como mecanismo de identificación}
Uno siempre tiene que ingresar una clave personal que debe ser obtenida de una manera determinada, y cuanto más complejo es el modo de entrecruzar información, más segura es la operatoria.
\end{example}

El sistema judicial también usa firma electrónica: los escritos judiciales son presentados virtualmente por los abogados con firma electrónica, y los jueces validan providencias proyectadas en el sistema de gestión a través de una clave alfanumérica de uso individual.

\subsection{Diferencia central con la firma digital y carga de la prueba}

Si el sistema de firma electrónica que se ha establecido en una aplicación es muy endeble, difícilmente pueda admitirse la emisión de una firma que no garantice identidad o que no garantice una plausible expresión de la voluntad.

\begin{table}[H]
\centering
\caption{Firma digital y firma electrónica}
\begin{tabularx}{\textwidth}{>{\raggedright\arraybackslash}p{0.30\textwidth} >{\raggedright\arraybackslash}X >{\raggedright\arraybackslash}X}
\toprule
 & \textbf{Firma digital} & \textbf{Firma electrónica} \\
\midrule
Entidad certificante & Sí -- organismo estatal & No \\
Plena fe & Sí -- presunción de autenticidad & No -- la parte que la presenta debe demostrar su seriedad \\
Carga de la prueba en caso de negativa & Quien impugna debe probar técnicamente la falsedad & Quien presenta debe probar la seriedad del sistema \\
Sistema de impugnación & Pericia informática que demuestre vicio en certificación & Cualquier medio para negar autenticidad \\
Uso típico & Trámites estatales, documentos oficiales & Homebanking, compras online, gestión judicial \\
\bottomrule
\end{tabularx}
\end{table}

\section{Documentos electrónicos sin firma}

Son el paralelo de los instrumentos particulares escritos. Incluyen todos los registros electrónicos —visuales, auditivos, videograbados— que permiten plasmar hechos o actos.

Con la proliferación de los teléfonos celulares inteligentes, videos y audios son elementos que están ``a flor de piel'' para documentar circunstancias, hechos y acontecimientos.

\begin{example}{Google Street View como registro electrónico}
En determinadas calles de la ciudad uno puede hacer la experiencia de situarse en un lugar determinado y ver cómo ese lugar fue cambiando con el paso del tiempo. Sin ir más lejos, el Paseo del Bajo en la Ciudad de Buenos Aires ha tenido una modificación significativa. Uno puede irse dos o tres años atrás y ver cómo ha evolucionado.
\end{example}

La aplicación Google Street View permite situarse en imágenes de múltiples lugares del mundo y ver la evolución de un lugar a lo largo del tiempo según las distintas ``pasadas'' de la cámara.

La validez de estos instrumentos electrónicos sin firma dependerá en cada caso de la calidad técnica que tengan, la posibilidad de escucha en audios, la calidad visual en videos, y de que existan otros elementos que permitan congruentemente valorarlos. En última instancia es función del juez valorar estos elementos en cada caso puntual.

\begin{table}[H]
\centering
\caption{Documentos electrónicos y su equivalente en papel}
\begin{tabularx}{\textwidth}{>{\raggedright\arraybackslash}p{0.28\textwidth} >{\raggedright\arraybackslash}p{0.28\textwidth} >{\raggedright\arraybackslash}X}
\toprule
\textbf{Documento electrónico} & \textbf{Equivalente en papel} & \textbf{Característica principal} \\
\midrule
Con firma digital & Instrumento público & Certificación estatal; plena fe de autenticidad \\
Con firma electrónica & Instrumento privado & Sistema de claves; carga de prueba en quien lo presenta \\
Sin firma & Instrumento particular & Videos, audios, registros; valoración judicial libre \\
\bottomrule
\end{tabularx}
\end{table}

% ============================================================
% CAPÍTULO 5
% ============================================================
\chapter{Correspondencia}

La correspondencia puede ser escrita en formato papel o electrónica (correo electrónico). Los artículos 318 y 319 del CCyC regulan su valor probatorio.

\begin{formula}{Artículos 318 y 319 CCyC — Correspondencia}
La correspondencia, cualquiera sea el medio empleado para crearla o transmitirla, puede presentarse como prueba por el destinatario, pero la que es confidencial no puede ser utilizada sin asentimiento del remitente. Los terceros no pueden valerse de la correspondencia sin asentimiento del destinatario, y del remitente si es confidencial.
\end{formula}

Reglas de la correspondencia como prueba:

\begin{enumerate}[label=\alph*)]
    \item El destinatario puede aportarla como prueba, siempre que no sea confidencial.
    \item Si es confidencial, requiere la conformidad del remitente para ser presentada en juicio. De lo contrario, se violaría la inviolabilidad de la correspondencia epistolar garantizada por el artículo 18 de la Constitución Nacional.
    \item Contenido comercial o negocial (tratativas para la celebración de un contrato, etc.): no habría inconveniente en incorporarla al proceso.
    \item El remitente puede presentarla; en ese caso, el destinatario debe reconocer o negar la recepción (no la confección).
    \item Correspondencia de terceros: la manera de autenticarla es a través de la prueba testimonial de los terceros a quienes se les atribuye la carta.
\end{enumerate}

En cuanto al valor probatorio: la correspondencia reconocida tiene valor de plena prueba, igual que el instrumento privado. Si es de un tercero que comparece como testigo y reconoce haber confeccionado la carta, vale como elemento indiciario según las circunstancias.

% ============================================================
% CAPÍTULO 6
% ============================================================
\chapter{Libros contables}

Antes de la reforma, los libros contables estaban regulados por el Código de Comercio. Hoy el Código Civil y Comercial, al haber unificado y derogado el Código de Comercio, tiene un capítulo dedicado a los libros contables.

Son un tipo especial de documentos privados de elaboración unilateral, pero que, pese a esa elaboración unilateral, su valor radica en que:

\begin{enumerate}
    \item Deben ser llevados con las formalidades exigidas por ley.
    \item Requieren de legalización por parte de la autoridad.
\end{enumerate}

\section{Valor probatorio de los libros contables (art. 330 CCyC)}

\begin{formula}{Artículo 330 CCyC — Valor probatorio de la contabilidad}
La contabilidad, obligada o voluntaria, llevada en la forma y con los requisitos prescritos, debe ser admitida en juicio como medio de prueba. Sus registros prueban contra quien la lleva. El juez tiene en todos los casos la facultad de apreciar esa prueba, y de exigir, si lo considera necesario, otra supletoria o complementaria.
\end{formula}

\begin{important}
Regla fundamental: la indivisibilidad. No se puede tener por válidos ciertos asientos que surjan de los libros de la contraparte y descartar otros. Si se admiten los libros, deben valorarse íntegramente. El artículo 330 establece que no pueden valorarse ciertos asientos y dejarse de lado otros.
\end{important}

\begin{example}{La indivisibilidad de los libros de comercio}
Si yo admito los asientos que surgen del libro de comercio debo admitirlos integralmente. Voy a admitir todo. No puedo servirme de unos y de los otros decir: no, esos no valen, los dejo de lado.
\end{example}

Casos de valoración de los libros contables:

\begin{table}[H]
\centering
\caption{Situaciones de valoración de los libros contables}
\begin{tabularx}{\textwidth}{>{\raggedright\arraybackslash}X >{\raggedright\arraybackslash}X}
\toprule
\textbf{Situación} & \textbf{Valor probatorio} \\
\midrule
Operaciones entre ambas partes que llevan contabilidad: los libros tienen asientos perjudiciales para quien los lleva & Plena prueba en contra de quien lleva los libros \\
Una parte presenta libros, la otra que debe llevarlos no los presenta & Los libros presentados se valoran favorablemente para quien los exhibe \\
Ambas partes presentan libros con contradicciones entre sí & El juez prescinde de los libros y se ciñe al resto de la prueba \\
Solo una parte lleva contabilidad (contraparte no comerciante) & Los libros valen como principio de prueba, evaluados con el resto de la prueba \\
\bottomrule
\end{tabularx}
\end{table}

Quien postula un desarrollo argumental en el proceso debe guardar coherencia con lo que surge de sus propios libros, porque todo aquello que surja de ellos en contra de lo que postula queda desvirtuado ante el juez.

% ============================================================
% CAPÍTULO 7
% ============================================================
\chapter{Ofrecimiento e impugnación de prueba documental}

\section{Oportunidad del ofrecimiento (arts. 333 y 335 CPCC)}

El momento adecuado y prácticamente exclusivo para la presentación de la prueba documental es:

\begin{itemize}
    \item \textbf{Con la demanda}: deben presentarse todos los documentos que la parte tiene en su poder. Los que no están en su poder deben individualizarse indicando dónde están.
    \item \textbf{Con la contestación de demanda}: el demandado deberá ofrecer su prueba documental en ese momento.
\end{itemize}

\begin{formula}{Artículo 333 CPCC — Documentos que deben acompañarse}
Con la demanda, reconvención y contestación de ambas, deberá acompañarse la prueba documental que estuviere en poder de las partes. Si no obrare en poder de las partes, estas la individualizarán, indicando su contenido, el lugar, archivo, oficina pública o persona en cuyo poder se encuentre.
\end{formula}

\subsection{Excepción: documentos posteriores (art. 335 CPCC)}

Si se obtiene un documento con posterioridad a la demanda, puede presentarse al proceso invocando el artículo 335, manifestando expresamente que ese documento llegó a poder de la parte con posterioridad a la demanda. Es conveniente acompañar alguna prueba demostrativa de esta circunstancia.

\begin{formula}{Artículo 335 CPCC — Documentos posteriores}
Después de interpuesta la demanda, no se admitirán al actor documentos, excepto los de fecha posterior, o aquellos en que el interesado jure o afirme no haber tenido antes conocimiento de ellos; en tales casos, se dará traslado a la otra parte, quien deberá cumplir la carga del artículo 356, inciso 1.
\end{formula}

\section{Documentos en poder de la contraparte (art. 388 CPCC)}

Cuando el documento está en poder de la otra parte, puede ofrecerse como prueba pidiendo al juez que intime a esa parte a presentarlos. Esta intimación tiene una característica peculiar: no se puede constreñir a la parte a presentar un documento que no se está seguro de que exista.

\begin{formula}{Artículo 388 CPCC — Documento en poder de una de las partes}
Si el documento se encontrare en poder de una de las partes, se le intimará a su presentación en el plazo que el juez determine. Cuando por otros elementos de juicio resultare manifiestamente verosímil su existencia y contenido, la negativa a presentarlo constituirá una presunción en su contra.
\end{formula}

\begin{example}{Los límites de la coacción sobre la parte}
Que vengan de las pestañas o que allanen domicilio para ir a buscar el documento de la casa es de película yanqui. No es de lo que permite el proceso civil. En materia civil, cuando algún documento no aparece, la herramienta que tiene la parte es la posibilidad de intimarlo bajo apercibimiento de que, si por otras circunstancias se llega a demostrar que es verosímil que el documento existía, la no presentación es una presunción en contra de quien omitió presentarlo.
\end{example}

Si a lo largo del proceso se ha demostrado que el contrato existía (porque hay otros elementos demostrativos de prestaciones cumplidas), el silencio o negativa frente a la intimación constituyen una presunción contraria al intimado.

\subsection{Documentos en poder de terceros}

Se puede pedir a los terceros que los presenten, pero el tercero puede oponerse si la presentación le puede generar un perjuicio. En ese caso el juez evaluará lo que corresponda.

Si el tercero presenta documentos, tiene derecho a pedir que se obtengan copias, que en las copias se certifique su autenticidad, y que se le devuelva el documento a la mayor brevedad.

\section{El traslado de la documentación}

De todo documento presentado en juicio se ordena el traslado a la otra parte. Dentro del plazo de ese traslado, la parte contraria está obligada a expedirse sobre la autenticidad de esos documentos.

\section{Impugnación del instrumento público (art. 395 CPCC)}

El mero desconocimiento de un instrumento público no alcanza para tacharlo de falso, dado que hace plena fe de su contenido. La impugnación requiere un procedimiento especial:

\begin{enumerate}
    \item \textbf{Paso 1}: dentro del plazo del traslado de la documentación, formular la impugnación: decir que ese instrumento público no es válido, se lo impugna.
    \item \textbf{Paso 2}: dentro de los 10 días hábiles judiciales contados desde el día siguiente de formulada la impugnación, promover el incidente de redargución de falsedad previsto por el artículo 395.
\end{enumerate}

\begin{formula}{Artículo 395 CPCC — Redargución de falsedad}
La redargución de falsedad de un instrumento público tramitará por incidente que deberá promoverse dentro del plazo de diez días de realizada la impugnación en el escrito en que se desconozca el instrumento, bajo apercibimiento de tenerlo por reconocido. Serán partes en el incidente quienes lo hubieren suscripto; si hubiere fallecido alguno de ellos, sus sucesores. El incidente tramitará en cuerda separada con el principal.
\end{formula}

El incidente de redargución de falsedad es un trámite incidental que corre en paralelo al proceso principal. Son parte necesaria en ese incidente:

\begin{itemize}
    \item El oficial público que intervino en la confección del instrumento.
    \item Las partes que intervinieron en el acto.
\end{itemize}

Contra todos ellos se dirige una acción de nulidad de la escritura. Este proceso de redargución de falsedad tramita en paralelo al juicio principal y, una vez que se encuentre en condiciones de resolver, se resuelve junto con la sentencia definitiva del juicio principal.

La tacha de falsedad también puede ejercerse por vía de acción civil o criminal (querella penal) de manera autónoma, no supeditada al proceso principal.

\subsection{Cuando NO se necesita la redargución de falsedad}

Cuando lo que se impugna es el contenido de las declaraciones de las partes ante el oficial público —y no la efectiva realización del acto, la fecha o el lugar—, no es necesaria la redargución de falsedad. En ese caso, cualquier prueba que se presente en el juicio principal puede desvirtuar los alcances de esas manifestaciones.

\begin{example}{Lo declarado ante el escribano}
Cuando lo que se está impugnando es lo que las partes dijeron al escribano el día de la escritura y que el escribano plasmó en el acta, ahí no sería necesaria la redargución.
\end{example}

% ============================================================
% CAPÍTULO 8
% ============================================================
\chapter{Prueba informativa}

\section{Concepto}

La prueba informativa es un modo de aportar al proceso hechos o actos registrados en documentos, archivos o registros que tiene algún tercero ajeno a las partes.

\begin{important}
La prueba informativa NO obtiene opiniones ni versiones sobre hechos (eso es la prueba testimonial). Lo que se busca es aportar datos registrados en archivos, documentos o registros que poseen terceros, sean entidades públicas o privadas.
\end{important}

\begin{example}{Oficio al SAME}
Oficio al SAME para que informe si obran en sus registros la atención médica brindada el día 14 de marzo de 2020 por un accidente en la vía pública ocurrido en Rivadavia al 5700, en el marco de un accidente de tránsito; a qué personas atendieron, o si atendieron a determinada persona y adónde la trasladaron.
\end{example}

\section{Formato del oficio (arts. 398 y 400 CPCC)}

La prueba informativa se materializa a través de un escrito denominado oficio, que tiene un formato determinado y debe contener:

\begin{enumerate}
    \item A quién va dirigido (por ejemplo, ``Al Director del Servicio de Emergencias SAME'').
    \item En qué actuaciones y ante qué juzgado tramitan.
    \item Domicilio del juzgado.
    \item Nombre del juez y del secretario, y un teléfono si es posible.
    \item Descripción detallada y minuciosa de lo que se requiere que informe.
    \item Transcripción de la orden dictada por el juez que autoriza ese pedido de informes.
    \item Mención de las sanciones aplicables por incumplimiento (artículo 398).
\end{enumerate}

\begin{formula}{Artículo 398 CPCC — Sanciones por incumplimiento}
Las entidades públicas o privadas que no contestaren el pedido de informes ordenado por el juez dentro del plazo de diez días serán pasibles de sanciones conminatorias progresivas que deberá establecer el juez en caso de incumplimiento.
\end{formula}

\begin{formula}{Artículo 400 CPCC — Oficios y firma}
Los oficios o exhortos que deban diligenciarse en el extranjero se ajustarán en cuanto a sus formas a lo establecido por los tratados o costumbre internacional. Los pedidos de informes podrán ser firmados directamente por el abogado de la parte, con la sola acreditación de su calidad de tal.
\end{formula}

\begin{example}{Modelo de oficio}
``Tengo el agrado de dirigirme a Usted, en los autos caratulados Fulano c/ Mengano sobre daños y perjuicios que tramitan ante el Juzgado Nacional de Primera Instancia en lo Civil N.\textdegree{} 18 a cargo del Dr. Fulanito, Secretaría N.\textdegree{} [tal] a cargo del Dr. Menganito, a efectos de requerirle tenga bien informar a este Juzgado acerca de la atención brindada a la Sra. Edith [Artal] en tal fecha en Rivadavia al 5700, detallando el diagnóstico que presentó y el lugar al que fue trasladada. La providencia o auto que ordena la presente medida dice: Buenos Aires, 15 de febrero del 2020: Líbrese el oficio solicitado por la parte actora. Firmado: Julián Domínguez, Juez Nacional en lo Civil.''
\end{example}

\section{Plazo y caducidad de la prueba informativa (art. 402 CPCC)}

La entidad que recibe el oficio tiene 10 días hábiles judiciales para contestar, contados desde el día siguiente a la presentación del oficio. El abogado debe acreditar en el expediente el diligenciamiento del oficio, mediante copia sellada como constancia de recepción.

\begin{formula}{Artículo 402 CPCC — Caducidad de la prueba informativa}
Si vencido el plazo fijado para contestar el informe la repartición pública o entidad privada no lo hubiere remitido, se tendrá por desistida de esa prueba a la parte que la pidió, sin sustanciación alguna, si dentro del quinto día no solicitare al juez la reiteración del oficio.
\end{formula}

\begin{table}[H]
\centering
\caption{Mecanismo de caducidad de la prueba informativa}
\begin{tabularx}{\textwidth}{>{\raggedright\arraybackslash}p{0.32\textwidth} >{\raggedright\arraybackslash}p{0.28\textwidth} >{\raggedright\arraybackslash}X}
\toprule
\textbf{Etapa} & \textbf{Plazo} & \textbf{Consecuencia del incumplimiento} \\
\midrule
Entidad debe contestar & 10 días hábiles desde la recepción del oficio & Puede operar la caducidad \\
Parte interesada debe pedir reiteración & 5 días desde el vencimiento de los 10 días & Si no pide reiteración: caducidad de la prueba informativa \\
Caducidad opera & Automáticamente al vencer los 5 días sin petición de reiteración & La parte pierde esa prueba definitivamente \\
\bottomrule
\end{tabularx}
\end{table}

\begin{important}
Es fundamental estar atentos, una vez obtenido el título y al ejercer la profesión, a que no se venzan los plazos de la prueba informativa: en ciertas circunstancias, su pérdida puede ser muy grave para los intereses del cliente.
\end{important}

\section{Artículo 38 CPCC: oficios dirigidos a altas autoridades}

El artículo 400 debe compatibilizarse con el artículo 38 del CPCC. Las comunicaciones por oficio dirigidas al Presidente de la Nación, ministros, secretarios de Estado, u otros jueces o magistrados (camaristas o ministros de la Corte) deben ser firmadas por el juez de la causa, no por el abogado.

En esos casos, el oficio tiene una redacción diferente y debe dejarse en secretaría para que sea firmado por el juez antes de diligenciarse.

\section{Modernización: sistemas electrónicos de pedido de informes}

Lentamente, la administración pública ha comenzado a darse cuenta de la obsolescencia del sistema de oficios en papel y ha incorporado sistemas electrónicos de comunicación oficial. Diversas reparticiones ya permiten que los pedidos de informes se hagan por vía electrónica:

\begin{itemize}
    \item Banco de la Nación Argentina.
    \item Banco Provincia de Buenos Aires.
    \item Banco Central de la República Argentina.
    \item Registro Nacional de la Propiedad Automotor.
    \item Registro de la Propiedad Inmueble.
    \item Registro Nacional de las Personas.
    \item Dirección Nacional de Migraciones.
    \item Varios bancos privados.
    \item Comisarías de la Ciudad de Buenos Aires y Policía Federal.
\end{itemize}

\begin{note}
El sistema tradicional de los oficios en formato papel —casi medievales— tiende a desaparecer. La cuarentena incrementó muchísimo el tráfico electrónico en materia procesal, cosa muy saludable que probablemente colabore para que cada vez en mayor medida se recurra a sistemas electrónicos.
\end{note}

Es importante conocer el sistema tradicional porque todavía está muy instalado en la idiosincrasia, y para pequeñas y medianas empresas y entidades privadas con poca comunicación con los tribunales, el oficio en papel sigue siendo la única forma de obtener la información requerida.

% ============================================================
% CAPÍTULO 9 — CORNELL
% ============================================================
\chapter{Repaso Cornell — Prueba documental e informativa}

El siguiente cuadro reúne, según el método Cornell, las preguntas centrales que permiten repasar activamente los contenidos desarrollados a lo largo de esta unidad, junto con su desarrollo correspondiente y una síntesis final.

\begin{longtable}{
    >{\raggedright\arraybackslash}p{0.30\textwidth}
    >{\raggedright\arraybackslash}p{0.62\textwidth}
}
\toprule
\textbf{Preguntas / palabras clave} & \textbf{Notas / respuestas} \\
\midrule
\endhead

\textbf{¿Qué es un documento en sentido amplio?} &
Todo elemento susceptible de representar una manifestación del pensamiento humano, independientemente de su soporte (papel, digital, magnético, en la nube). El concepto fue ampliado significativamente por el CCyC 2015. \\

\textbf{¿Cuáles son las categorías principales de documentos?} &
Soporte papel (instrumento público, instrumento privado, instrumento particular sin firma) y soporte electrónico (firma digital, firma electrónica, sin firma). Además: correspondencia y libros contables. \\

\textbf{¿Qué hace ``plena fe'' en un instrumento público?} &
La celebración del acto, la fecha, el lugar y los hechos que el oficial público enuncia como cumplidos ante él (art. 296 CCyC). NO hace plena fe el contenido de las declaraciones de las partes, que admite prueba en contrario. \\

\textbf{¿Cómo se impugna un instrumento público?} &
Se formula la impugnación al correr el traslado y, dentro de los 10 días hábiles siguientes, se promueve el incidente de redargución de falsedad (art. 395 CPCC). El oficial público es parte necesaria. También puede ejercerse por acción autónoma civil o penal. \\

\textbf{¿Cuándo NO es necesaria la redargución?} &
Cuando lo impugnado son las declaraciones de las partes ante el escribano (no la efectiva realización del acto). En ese caso, cualquier medio de prueba en el juicio principal alcanza. \\

\textbf{¿Qué es el cotejo caligráfico y cuándo se usa?} &
Es la pericia que compara la firma dubitada (en duda) con firmas indubitadas (auténticas) de la persona que niega haberlo firmado. Se usa cuando el demandado niega su firma en un instrumento privado. Se realiza mediante documentos indubitados acordados por las partes o un cuerpo de escritura confeccionado en presencia del perito (arts. 390/458 CPCC). \\

\textbf{¿Qué es el abuso de firma en blanco?} &
Cuando alguien suscribe una hoja en blanco por confianza y la otra parte la llena con contenido distinto al convenido. Para invocarlo se requiere un principio de prueba por escrito (art. 315 CCyC). No bastan solos los testimonios. \\

\textbf{¿Cuál es la diferencia entre firma digital y firma electrónica?} &
La firma digital tiene entidad estatal certificante, hace plena fe y quien la impugna debe probar el vicio técnico (Ley 25.506). La firma electrónica (homebanking, compras online) no tiene certificante; quien la presenta tiene la carga de demostrar la seriedad del sistema. \\

\textbf{¿Cuándo debe ofrecerse la prueba documental?} &
Con la demanda (actor) o con la contestación de demanda (demandado). Solo pueden presentarse después documentos de fecha posterior o cuyo desconocimiento se justifique bajo juramento (art. 335 CPCC). \\

\textbf{¿Qué pasa si la contraparte intimada no presenta el documento?} &
Si a lo largo del proceso se demuestra que era verosímil la existencia de ese documento, la negativa constituye una presunción contraria (art. 388 CPCC). No hay allanamiento ni coacción física; es solo una presunción en contra. \\

\textbf{¿Qué es la prueba informativa?} &
Es el pedido de informes a entidades públicas o privadas para que informen datos que figuran en sus registros. No obtiene opiniones (eso es testimonial), sino datos registrados. Se concreta a través de un oficio judicial firmado por el abogado (art. 400 CPCC). \\

\textbf{¿Cómo opera la caducidad de la prueba informativa?} &
La entidad tiene 10 días hábiles para contestar. Si no contesta, la parte interesada tiene 5 días para pedir la reiteración del oficio. Si no lo pide, opera la caducidad: pierde definitivamente esa prueba (art. 402 CPCC). \\

\textbf{¿Qué valor tienen los libros contables?} &
Prueban plenamente en contra de quien los lleva (asientos perjudiciales). Si la contraparte no presenta los suyos, los presentados valen favorablemente. Si ambas partes los presentan y son contradictorios, el juez prescinde de ellos. Son indivisibles: no se puede tomar solo lo favorable (art. 330 CCyC). \\

\textbf{¿Qué permite el art. 18 CN respecto de la correspondencia?} &
La inviolabilidad de la correspondencia epistolar. Por eso la correspondencia confidencial no puede ser presentada en juicio sin el consentimiento del remitente. La correspondencia de contenido comercial o negocial sí puede presentarse sin esa restricción. \\

\midrule
\multicolumn{2}{p{0.94\textwidth}}{
\textbf{Resumen:}
Los medios de prueba documentales en el proceso civil argentino abarcan desde el ticket de estacionamiento hasta la firma digital certificada por el Estado, en virtud del concepto amplio de documento incorporado por el CCyC 2015. El valor probatorio de cada categoría varía: los instrumentos públicos hacen plena fe de la realización del acto, fecha y lo actuado ante el oficial público, pero admiten prueba en contrario respecto de las declaraciones de las partes; los instrumentos privados valen plenamente cuando son reconocidos o autenticados por cotejo caligráfico; y los instrumentos particulares son valorados libremente por el juez según la congruencia del caso. En el plano electrónico, la firma digital (Ley 25.506) otorga plena fe, mientras que la firma electrónica —más frágil— impone la carga de la prueba a quien la presenta. La prueba informativa, por su parte, permite incorporar al proceso datos objetivamente registrados en archivos de terceros mediante oficios con plazo de 10 días y caducidad a los 5 días de vencido ese plazo si no se reitera. El correcto manejo de estos medios, sus plazos y sus mecanismos de impugnación resulta decisivo para el éxito o fracaso de una pretensión judicial.
} \\

\bottomrule
\end{longtable}

\end{document}
